% =====================================================================
% TLECUIL Software Quality Assurance and Verification & Validation Plan
% Build:  latexmk -pdf tlecuil_vv_plan.tex   (same folder as references.bib)
% =====================================================================
\documentclass[11pt]{article}
\usepackage[margin=1in]{geometry}
\usepackage[T1]{fontenc}
\usepackage{lmodern}
\usepackage{amsmath}
\usepackage{booktabs,longtable,array,tabularx}
\usepackage{xcolor}
\usepackage[most]{tcolorbox}
\usepackage{enumitem}
\usepackage[style=ieee,backend=biber,doi=true,url=true]{biblatex}
\addbibresource{references.bib}
\usepackage[colorlinks=true,linkcolor=blue!50!black,citecolor=blue!50!black,urlcolor=blue!50!black]{hyperref}

\newtcolorbox{plain}{colback=violet!5,colframe=violet!60!black,title=In plain words,fonttitle=\bfseries\small}
\newtcolorbox{decision}{colback=orange!6,colframe=orange!60!black,title=Decision needed,fonttitle=\bfseries\small}
\newcommand{\tlecuil}{\textsc{Tlecuil}}
\newcolumntype{L}{>{\raggedright\arraybackslash}X}

\title{\tlecuil{} Software Quality Assurance and\\Verification \& Validation Plan}
\author{Ira Strong}
\date{Revision 0 (draft) --- \today\\[4pt]\small Applies to \tlecuil{} v0.2.x}

\begin{document}
\maketitle

\begin{abstract}
This plan defines how \tlecuil{}, an open first-order fission surface power sizer, will be
verified, validated and controlled so that it can later be qualified for use in a formal
quality assurance program. It classifies the software under ASME NQA-1-2024, DOE O 414.1F
and NASA NPR 7150.2D. It then states the software requirements and traces each one to an
equation, a function and a test. Finally, it sets out the verification and validation
approach, configuration management, problem reporting, reviews and records. It is
written to be usable by students as an introduction to software quality assurance for
engineering codes.
\end{abstract}

\tableofcontents

%=====================================================================
\section{Purpose and scope}
%=====================================================================
\tlecuil{} computes a steady-state thermal--electrical balance of a fission surface power
system: reactor thermal power, conversion efficiency, waste-heat loads, radiator areas and
U-235 burned. Its equations are documented in \emph{TLECUIL: A First-Order Fission Surface
Power Sizer} (the ``Theory'' document).

This plan covers the Python engine (\texttt{tlecuil.py}), the benchmark and test
scripts, the JavaScript copy of the engine in the interactive page, and the
documentation. It does not cover any future mass, shielding or efficiency-map modules;
each will be added to this plan when it is developed.

\begin{plain}
Software quality assurance (SQA) is a set of habits that let someone else trust your
code without re-deriving it. Write down what the code must do; show that it does it;
control every change; record every error and how it was fixed; and have someone
independent check the work.
\end{plain}

%=====================================================================
\section{What kind of software is \tlecuil{}?}\label{sec:class}
%=====================================================================
The answer depends on who uses it and for what. The same code can be ``research
software'' in a thesis and ``design analysis software'' inside a licensed project. The
grade is set by the \emph{use}, not by the code itself.

\subsection{ASME NQA-1-2024, Part II, Subpart 2.7}
NQA-1 is the consensus quality assurance standard for nuclear facility applications.
Subpart 2.7 sets requirements for computer software used in that work
\cite{nqa1_2024}. Within an NQA-1 program, \tlecuil{} would be \emph{custom-developed
design and analysis software}: a program whose results support design outputs.

Today \tlecuil{} is developed outside any organization's NQA-1 program, so NQA-1 does not
formally apply. If an organization later wants to use it for design work, that
organization would have to bring it under its own program, either by qualifying it as
acquired or existing software or by re-developing it under its program. The aim of
this plan is to make that a review of existing records rather than a rework.

\subsection{DOE O 414.1F}
DOE O 414.1F, approved in August 2026, defines safety software as including, at a
minimum, software cited in a DOE-approved safety analysis or approved hazard analysis.
It also requires appropriate software engineering consensus standards, without naming
NQA-1 specifically \cite{doe4141f}. \tlecuil{} is \textbf{not safety software} under this
definition. It would become safety software only if a DOE contractor cited it in such an
analysis. The contractor's quality assurance program would then grade it.

\subsection{NASA NPR 7150.2D}
NASA classifies software from Class A (human-rated) to Class F (business and IT)
\cite{npr7150_2d}.
\begin{itemize}
  \item Class E covers design-concept, research, technology and general-purpose software.
        This describes \tlecuil{} as a research and teaching tool \textbf{today}.
  \item Class D covers basic science and engineering design, and research and technology
        software. \tlecuil{} would move to Class D if NASA used it to inform engineering
        design decisions.
  \item It is not safety-critical: it does not control hardware or protect people.
\end{itemize}

\subsection{Classification used in this plan}
\begin{center}
\begin{tabularx}{\textwidth}{lL}
\toprule
Attribute & Determination \\
\midrule
Software type & Custom-developed scientific / engineering analysis software \\
Current use & Research, education, cross-checking published first-order results \\
Safety software (DOE O 414.1F) & No, unless cited in a DOE-approved safety or hazard analysis \\
NASA class (NPR 7150.2D) & E now; D if used for NASA engineering design \\
NQA-1 Subpart 2.7 & Not invoked (no NQA-1 program); developed to be qualifiable as design and analysis software \\
Grading of this plan & ``NQA-1-ready'' graded approach: the controls in Sections~\ref{sec:req}--\ref{sec:records}, with records kept so a later qualification can rely on them \\
\bottomrule
\end{tabularx}
\end{center}

\begin{decision}
Before \tlecuil{} is used in any project under an employer's or sponsor's quality assurance
program, confirm ownership and approvals for the code, and have that program's software
quality authority confirm the classification above.
\end{decision}

%=====================================================================
\section{Verification and validation framework}\label{sec:framework}
%=====================================================================
The plan follows the terminology of Oberkampf and Roy \cite{oberkampf2010vv}, the
validation-uncertainty approach of ASME V\&V~20 \cite{asme_vv20}, and the life-cycle V\&V
structure of IEEE 1012 \cite{ieee1012_2024}.
\begin{description}[style=nextline]
  \item[Code verification] \emph{Is the math coded correctly?} The code is compared with
        exact answers: hand calculations, closed-form limits, conservation laws and
        independent numerical integration.
  \item[Solution verification] \emph{How large are the numerical errors in a given run?} For
        \tlecuil{} these come from the series truncation in the radiator integral and from
        the optimiser's tolerance. Both are bounded and tested.
  \item[Validation] \emph{Does the model describe reality well enough for its intended
        use?} The model is compared with measurements from hardware, with the experimental
        and input uncertainties quantified. The validation comparison error is $E=S-D$
        (simulation minus data), judged against the validation uncertainty $u_\text{val}$
        \cite{asme_vv20}.
  \item[Benchmarking] Comparison with other models or published design studies. It is
        useful, but it is not validation (Theory, Sec.~11).
\end{description}

%=====================================================================
\section{Software requirements}\label{sec:req}
%=====================================================================
Each requirement has an identifier so that it can be traced. ``Shall'' marks a binding
requirement.
\begin{longtable}{>{\ttfamily}p{0.09\textwidth}p{0.83\textwidth}}
\toprule
ID & Requirement \\
\midrule
\endhead
FR-01 & The code shall compute converter gross power from delivered power, PMAD efficiency, and proportional and fixed parasitic loads (Theory Eq.~3.2). \\
FR-02 & The code shall compute conversion efficiency as a fixed value or as a fraction of Carnot (Theory Eq.~4.4). \\
FR-03 & The code shall reject any efficiency above the Carnot limit (Theory Eq.~4.3). \\
FR-04 & The code shall compute reactor thermal power including reactor and piping losses (Theory Eq.~4.5). \\
FR-05 & The code shall compute main radiator radiating area for isothermal and pumped-loop radiators, including panel effectiveness and margin (Theory Eqs.~5.5--5.9). \\
FR-06 & The code shall report panel planform as radiating area divided by the number of radiating sides (Theory Eq.~5.2). \\
FR-07 & The code shall compute the electronics radiator area (Theory Sec.~6). \\
FR-08 & The code shall compute the U-235 mass fissioned (Theory Eq.~7.1). \\
FR-09 & The code shall find the cold-end temperature that minimises radiator area (Theory Sec.~8). \\
FR-10 & The code shall report the energy-balance residual on every run (Theory Eq.~6.2). \\
FR-11 & The code shall reject non-physical inputs with an explanatory message. \\
FR-12 & The code shall impose no upper limit on delivered power. \\
NR-01 & Results shall agree with exact solutions to a relative error $\le10^{-9}$. \\
NR-02 & The radiator integral shall be accurate to $\le10^{-9}$ relative error for sink temperatures from $10^{-9}$\,K to within 1\,\% of the outlet temperature. \\
IR-01 & The JavaScript engine shall reproduce the Python engine to a relative difference $\le10^{-9}$ on all parity cases. \\
DR-01 & Every equation shall be derived and cited in the Theory document, and every code reference to it shall carry its LaTeX label and current number. \\
DR-02 & Every published input used in a benchmark shall be traceable to its source, with any interpretation disclosed. \\
\bottomrule
\end{longtable}

%=====================================================================
\section{Requirements traceability matrix}\label{sec:rtm}
%=====================================================================
{\small\begin{longtable}{>{\ttfamily}p{0.09\textwidth}p{0.2\textwidth}p{0.25\textwidth}p{0.36\textwidth}}
\toprule
Req. & Theory & Code & Test(s) in \texttt{test\_\allowbreak{}tlecuil.py} \\
\midrule
\endhead
FR-01 & Eqs.~3.1--3.3 & \texttt{size\_\allowbreak{}fsp} & \texttt{hand\_\allowbreak{}calculation}; \texttt{energy\_\allowbreak{}conserved\_\allowbreak{}everywhere}; benchmark S1--S3 \\
FR-02 & Eq.~4.4 & \texttt{conversion\_\allowbreak{}efficiency} & \texttt{optimum\_\allowbreak{}matches\_\allowbreak{}closed\_\allowbreak{}form} \\
FR-03 & Eq.~4.3 & \texttt{validate} & \texttt{second\_\allowbreak{}law\_\allowbreak{}guard} \\
FR-04 & Eq.~4.5 & \texttt{size\_\allowbreak{}fsp} & \texttt{energy\_\allowbreak{}conserved\_\allowbreak{}everywhere} \\
FR-05 & Eqs.~5.5--5.9 & \texttt{radiating\_\allowbreak{}area\_\allowbreak{}m2}, \texttt{\_\allowbreak{}integral\_\allowbreak{}inv\_\allowbreak{}t4} & \texttt{fluid\_\allowbreak{}loop\_\allowbreak{}matches\_\allowbreak{}quadrature}; \texttt{fluid\_\allowbreak{}loop\_\allowbreak{}limit\_\allowbreak{}is\_\allowbreak{}isothermal}; \texttt{zero\_\allowbreak{}kelvin\_\allowbreak{}sink\_\allowbreak{}branch}; \texttt{small\_\allowbreak{}and\_\allowbreak{}large\_\allowbreak{}sinks\_\allowbreak{}match\_\allowbreak{}quadrature}; benchmarks B1--B4 \\
FR-06 & Eq.~5.2 & \texttt{size\_\allowbreak{}fsp} & \texttt{two\_\allowbreak{}sided\_\allowbreak{}panel\_\allowbreak{}halves\_\allowbreak{}planform\_\allowbreak{}not\_\allowbreak{}radiating\_\allowbreak{}area} \\
FR-07 & Sec.~6 & \texttt{size\_\allowbreak{}fsp} & electronics comparison (information only) \\
FR-08 & Eq.~7.1 & \texttt{size\_\allowbreak{}fsp} & \texttt{fuel\_\allowbreak{}rule\_\allowbreak{}of\_\allowbreak{}thumb} \\
FR-09 & Sec.~8 & \texttt{optimal\_\allowbreak{}cold\_\allowbreak{}temperature} & \texttt{optimum\_\allowbreak{}matches\_\allowbreak{}closed\_\allowbreak{}form} \\
FR-10 & Eq.~6.2 & \texttt{size\_\allowbreak{}fsp} & \texttt{energy\_\allowbreak{}conserved\_\allowbreak{}everywhere} \\
FR-11 & -- & \texttt{validate} & \texttt{second\_\allowbreak{}law\_\allowbreak{}guard} (planned: one test per rule, VT-07) \\
FR-12 & Sec.~9 & all & \texttt{linear\_\allowbreak{}scaling\_\allowbreak{}without\_\allowbreak{}fixed\_\allowbreak{}loads} (to 5000\,kWe) \\
NR-01 & -- & all & all verification tests \\
NR-02 & Eqs.~5.8--5.9 & \texttt{\_\allowbreak{}integral\_\allowbreak{}inv\_\allowbreak{}t4} & \texttt{small\_\allowbreak{}and\_\allowbreak{}large\_\allowbreak{}sinks\_\allowbreak{}match\_\allowbreak{}quadrature} \\
IR-01 & -- & \texttt{web/tlecuil.html} & \texttt{check\_\allowbreak{}parity.js} (CI) \\
DR-01 & all & comments & \texttt{docs/sync\_\allowbreak{}refs.py --check} (CI) \\
DR-02 & Sec.~11 & \texttt{benchmarks.py} & benchmark tests; independent review IR-001 \\
\bottomrule
\end{longtable}}

%=====================================================================
\section{Verification test plan}\label{sec:vtest}
%=====================================================================
\subsection{Existing tests (v0.2)}
Sixteen automated tests run on every change. They fall into four kinds:
\begin{enumerate}
  \item \textbf{Hand calculations}: one complete case checked to $10^{-12}$.
  \item \textbf{Exact limits}: the fluid-loop radiator reduces to the isothermal one; the
        optimiser reproduces the closed-form optimum; the $T_s=0$ branch is exact.
  \item \textbf{Independent numerics}: the radiator integral against SciPy adaptive
        quadrature over ten sink temperatures.
  \item \textbf{Conservation and invariants}: the energy-balance residual, linear scaling,
        and the two-sided-panel rule (a regression test for problem report PR-001).
\end{enumerate}

\subsection{Planned additions}
\begin{longtable}{>{\ttfamily}p{0.08\textwidth}p{0.84\textwidth}}
\toprule
ID & Planned verification activity \\
\midrule
\endhead
VT-01 & Property-based tests (random valid inputs): energy balance closes; area falls as emissivity, effectiveness or radiator temperature rises; area scales linearly with heat load. \\
VT-02 & Golden-file regression: freeze the outputs of all parity cases at each release; any change must be explained in the change log. \\
VT-03 & Optimiser solution verification: show the optimum moves by less than its tolerance when the tolerance is halved. \\
VT-04 & Series-truncation check: the 20-term series against 40 terms and quadrature across $a/T_\text{out}\in[0,0.5)$. \\
VT-05 & Parity of the JavaScript engine extended to every preset and to error paths (same inputs rejected, same messages). \\
VT-06 & Independent hand verification of the worked example by a second person, recorded with name and date. \\
VT-07 & One test per input-validation rule (FR-11). \\
\bottomrule
\end{longtable}

\subsection{Acceptance criteria}
All tests pass on the supported Python versions; the parity check reports a maximum
relative difference $\le10^{-9}$; \texttt{sync\_refs.py --check} passes; and the Theory PDF
builds without undefined references.

%=====================================================================
\section{Validation plan}\label{sec:val}
%=====================================================================
\subsection{Intended use and required accuracy}
Validation is always \emph{for} an intended use. For first-order trade studies, the
proposed target is: reactor thermal power within $\pm5\,\%$ and radiator area within
$\pm15\,\%$ of what measured hardware would need, with the uncertainty stated. These
targets are proposals and must be confirmed by the intended user.

\subsection{Validation hierarchy}
A complete system has never been built, so validation proceeds from the bottom up:
\begin{enumerate}
  \item \textbf{Component level, power conversion.} Compare the efficiency model with
        measured converter data, for example the Technology Demonstration Unit tests
        \cite{briggs2016tdu}.
  \item \textbf{Component level, heat rejection.} Compare the radiator model, including
        panel effectiveness, with thermal-vacuum test data from radiator panels. Such data
        still has to be identified. This is the most important gap, because the panel
        effectiveness is currently calibrated on a design study, not on hardware.
  \item \textbf{Integrated non-nuclear system.} Compare whole-system energy flows with
        electrically heated system tests such as the TDU.
  \item \textbf{System level.} Only design studies exist; these remain benchmarks.
\end{enumerate}

\subsection{First validation case (VC-01): TDU energy flows}
The TDU tests report converter power, converter efficiency, net system power and system
efficiency at nominal and maximum conditions \cite{briggs2016tdu}. A preliminary check at
the nominal point shows the case is useful:
\begin{itemize}
  \item Converter: 8.9\,kW at 21.7\,\% implies a converter heat input of $8.9/0.217=41.0$\,kW.
  \item Net 8.1\,kW at a system efficiency of 17.2\,\% implies a total heat input of
        $8.1/0.172=47.1$\,kW.
  \item The difference, about 13\,\% of the total, is heat that never reached the converter.
        That is exactly what \tlecuil{}'s reactor heat-loss fraction (Theory Eq.~4.5)
        represents.
\end{itemize}
This is a consistency check, not yet a validation. VC-01 will document the TDU test
conditions, the measurement uncertainties from the source, and an a-priori \tlecuil{}
prediction made \emph{before} comparing, and then report $E$ and $u_\text{val}$.

%=====================================================================
\section{Configuration management}\label{sec:cm}
%=====================================================================
\begin{itemize}
  \item \textbf{Repository.} All code, tests, documents and records live in one Git
        repository. The \texttt{main} branch is protected: changes arrive by pull
        request, and each pull request must pass CI.
  \item \textbf{Versions.} Semantic versioning (MAJOR.MINOR.PATCH). A change that alters any
        numerical result increments at least MINOR and is described in
        \texttt{CHANGELOG.md}.
  \item \textbf{Baselines.} Each release is a signed Git tag. Its record lists the commit,
        the Python and package versions, the test results, and the PDF built from that
        commit.
  \item \textbf{Environment.} Supported Python versions and package versions are listed in
        \texttt{requirements.txt} and exercised by CI.
  \item \textbf{Generated artefacts.} Benchmark tables and figures are generated by the
        code and are never edited by hand.
\end{itemize}

%=====================================================================
\section{Problem reporting and corrective action}\label{sec:prca}
%=====================================================================
Every error found after a result has been shared is recorded as a problem report (PR)
using the repository's issue template. Each report records: what is wrong; which versions
and results are affected; its impact (does it make past results non-conservative?); the
cause; the fix; the regression test that now guards it; and whether users who received
affected results were told.

\begin{longtable}{>{\ttfamily}p{0.08\textwidth}p{0.42\textwidth}p{0.4\textwidth}}
\toprule
ID & Problem & Disposition \\
\midrule
\endhead
PR-001 & v0.1 divided by the number of radiating sides twice; its reported ``emitting area'' was half the true radiating area for two-sided panels. Non-conservative. & Fixed in v0.2. Regression test \texttt{two\_sided\_\allowbreak{}panel\_\allowbreak{}\ldots}. User informed (revision history in Theory). \\
PR-002 & The closed-form radiator integral lost precision for sink temperatures near 0\,K and could return a negative area at $10^{-3}$\,K. Found by independent review. & Fixed in v0.2 (series branch). Test \texttt{small\_and\_\allowbreak{}large\_\allowbreak{}sinks\_\allowbreak{}\ldots}. No published result affected (all used $T_s\ge185$\,K). \\
\bottomrule
\end{longtable}

%=====================================================================
\section{Reviews}\label{sec:reviews}
%=====================================================================
\begin{itemize}
  \item \textbf{Independent technical review} of each release that changes the physics,
        by someone who did not write the change. The record keeps the reviewer, date,
        scope, findings and their disposition.
  \item \textbf{IR-001 (v0.2).} Independent review of the equations, the code--document
        consistency, the benchmark inputs against primary sources, and the citations.
        Result: no critical findings. Four major findings (overstated independence of the
        benchmarks, an undisclosed interpretation of a 2\,kWe load, PR-002, and an
        unsupported claim) and fifteen minor findings, all dispositioned in v0.2.
\end{itemize}

%=====================================================================
\section{Documentation and records}\label{sec:records}
%=====================================================================
\begin{center}
\begin{tabularx}{\textwidth}{lL}
\toprule
Document & Where \\
\midrule
SQA and V\&V plan & This document \\
Software requirements & Section~\ref{sec:req} \\
Design and theory & \texttt{docs/tlecuil\_theory.tex} \\
Traceability & Section~\ref{sec:rtm}; \texttt{sync\_refs.py} \\
Test results & CI logs per commit; release test record \\
Benchmark report & Theory Sec.~11; \texttt{benchmark\_results.json} \\
Validation reports & To be written per case (VC-01 first) \\
User guide & \texttt{README.md}; Theory worked example \\
Problem reports, reviews & Repository issues; Sections~\ref{sec:prca}--\ref{sec:reviews} \\
Change history & \texttt{CHANGELOG.md}; Git history \\
\bottomrule
\end{tabularx}
\end{center}
Records are kept for the life of the software and are never rewritten. A correction is
made by a new entry that refers to the old one.

%=====================================================================
\section{Roadmap to a qualifiable baseline}
%=====================================================================
\begin{enumerate}
  \item \textbf{v0.2 (now):} requirements, traceability, 16 tests, CI, problem reports,
        one independent review.
  \item \textbf{v0.3:} VT-01 to VT-07; golden-file regression; first signed release record.
  \item \textbf{v0.4:} VC-01 validation report (TDU); search for radiator-panel test data.
  \item \textbf{v1.0:} a frozen baseline with a complete record set, ready for a quality
        program to review for acceptance. New modules (mass, shielding, efficiency maps)
        enter only with their own requirements, tests and validation cases.
\end{enumerate}

\setcounter{biburlnumpenalty}{9000}\setcounter{biburllcpenalty}{7000}\setcounter{biburlucpenalty}{8000}
\printbibliography[heading=bibintoc]
\end{document}
